\selectlanguage{spanish}
\Language{es}{Español}{Guía de uso}
\firsttopic{Abrir, suspender y volver}
Inkline ejecuta una shell local en reMarkable 2. Use un Type Folio o un teclado
USB, trabaje en la tableta o conéctese mediante SSH o Goblin Mosh.
El texto y los gráficos se muestran en escala de grises.

\begin{notice}
\key{Opt+RightAlt+T} abre Inkline desde la propia tableta.\par
\key{Botón de encendido:} pulse y suelte para suspender; vuelva a pulsar para
reanudar. Las shells y los búferes del editor permanecen en RAM.
\end{notice}
La suspensión pausa la tableta sin cerrar Inkline ni guardar ajustes en la
memoria flash. Las conexiones de red pueden necesitar reconectarse. La pulsación
que despierta la tableta no vuelve a suspenderla. Después de actualizar,
cierre y abra Inkline una vez para activar el nuevo manejo del botón.

Para volver a los cuadernos, toque \key{Quit} o pulse \key{Ctrl+Shift+Q} y confirme.
\code{exit} cierra la terminal actual; cerrar la última devuelve los cuadernos.
La barra inferior contiene Escape, historial arriba/abajo, Quit y Settings.
\code{clear} borra la breve ayuda inicial y el pequeño logotipo Goblin.

Si una aplicación lanzada por atajo ocupa la pantalla, \key{Opt+RightAlt+T} vuelve a Inkline.
En caso de emergencia, mantenga \key{Opt+RightAlt+Backspace durante dos segundos}.
La recuperación cierra todas las terminales y pierde el trabajo sin guardar
de esas sesiones. Al reiniciar la tableta, los cuadernos siguen siendo la interfaz inicial.

RightAlt es la tecla Alt/Opt a la derecha de la barra espaciadora. Manténgala junto con la tecla Opt independiente y pulse T, sin Ctrl. Super solo sigue disponible para sus propios atajos.

\ProjectLinks

\topic{Teclas y posiciones de terminal}
El Type Folio tiene una tecla \key{Opt} independiente entre Ctrl y Alt, y una
tecla \key{Alt/Opt derecha}. Cumplen funciones distintas.

\begin{keytable}
\key{Teclas} & \key{Acción} \\ \midrule
Opt independiente + 1--0 & F1--F10 \\
Alt/Opt derecha + 1--9 & Elegir posición de terminal 1--9 \\
Alt/Opt derecha + Left / Right & Posición anterior / siguiente \\
Alt/Opt derecha + Tab & Escape \\
Alt/Opt derecha + Up / Down & PageUp / PageDown \\
Alt/Opt derecha + Backspace & Confirmar el cierre de todas las terminales \\
Cualquier Alt + Space & Teclado Unicode \\
Alt/Opt derecha + C / V & Copiar selección / pegar \\
Ctrl+Shift+B & Ocultar / mostrar barra inferior \\
\end{keytable}

En el \key{Type Folio estadounidense}, Alt/Opt derecha+\key{0} escribe \code{+}.
Alt/Opt derecha con la \key{tecla menos justo a la izquierda de Backspace}
escribe \code{=}. Con el método de entrada desactivado, Shift+6 introduce un
circunflejo literal: \code{c^2} sigue siendo \code{c^2}. No hay atajos de zoom.

Los teclados USB conservan sus teclas de función habituales. El atajo de Opt
independiente usa Meta cuando el teclado USB lo proporciona.

Hay hasta \key{nueve posiciones independientes}. Elegir una vacía abre una shell.
El indicador cuenta las terminales abiertas: si solo están abiertas 1 y 9,
la 9 muestra \key{Terminal 2 / 2 · slot 9}. Cerrar una shell libera su posición
sin cambiar la numeración de los atajos.

Cada terminal tiene su shell, directorio, composición de entrada, gráficos y
500 líneas de historial. Los ajustes de pantalla y método de entrada son comunes.

\topic{Ajustes y pantalla}
Toque \key{Settings}, el quinto botón inferior. Si la barra está oculta, pulse
\key{Alt+Space}, después \key{F2}, o toque Settings en el teclado Unicode.
La opción rellena está activa; el borde discontinuo indica el foco del teclado.
Tab mueve el foco; las flechas y Enter cambian la selección.

\key{Caps Lock:} elija Control (predeterminado) o Caps Lock.
\key{Tamaño:} pellizque con dos dedos o use los botones menos/más de Settings.
Los movimientos lentos ajustan de un píxel en un píxel, entre \key{6 y 48 píxeles}.

\key{Text darkness:} 50\% es el renderizado original. Los valores superiores
oscurecen los bordes. \key{Minimum contrast} separa la luminancia de texto y
fondo cuando una aplicación escoge colores parecidos; su valor inicial es 35\%.

\begin{choices}
\key{Modo e-paper} & \key{Compromiso} \\ \midrule
Fast & Espera mínima; forma de onda de animación; más imágenes residuales \\
Balanced & Forma de onda de interfaz; agrupa salida durante 32 ms \\
Crisp (inicial) & Grises más limpios; forma de onda de contenido; más lento \\
Mono & Blanco y negro rápido; descarta grises de las imágenes \\
Saver & Agrupa durante 120 ms; menos actualizaciones en ráfagas \\
\end{choices}

Solo se redibujan las filas modificadas, pero el panel sigue necesitando tiempo.
No se ha medido el ahorro de batería de Saver. Actualizar conserva el modo
guardado explícitamente, sin imponer el nuevo valor predeterminado.

\begin{notice}
Todos los ajustes permanecen en \key{RAM} y se guardan juntos al salir normalmente
de Inkline. Levantar el dedo, cerrar Settings o suspender no escribe en flash.
Una sesión sin cambios no escribe ajustes. Un fallo, cierre forzado o corte de
energía puede perder los cambios que aún no se hayan guardado.
\end{notice}

\topic{Unicode y métodos de entrada}
\key{Alt+Space} abre el teclado Unicode. Toque un carácter para insertarlo;
el teclado permanece abierto. Toque \key{Done} o pulse Escape para cerrarlo.

Las categorías incluyen letras, marcas, números, puntuación, matemáticas/monedas,
símbolos, espacios, caracteres de formato y uso privado. Arrastre verticalmente
o use flechas, PageUp/PageDown, Home/End y la rueda del ratón.
\key{Tab} cambia de categoría; Enter inserta el carácter seleccionado.

Escriba un valor Unicode hexadecimal, como \code{03bb}, y Enter para insertar λ.
Backspace edita el valor. El círculo punteado de las marcas combinantes es una
ayuda visual; solo se inserta la marca. Los espacios y caracteres de formato
tienen etiquetas. El catálogo usa las tablas Unicode de Qt y excluye controles,
valores sin asignar, sustitutos y no caracteres. Algunos glifos necesitan fuentes
no instaladas; los valores de uso privado requieren una fuente adecuada.

En el selector, pulse \key{F2} o toque \key{Settings} y elija \key{Input method}:
\begin{choices}
Off — direct keyboard & Escritura estadounidense normal \\
Romaji & Hiragana japonés desde entrada romanizada \\
Pinyin & Chino mediante pinyin \\
Zhuyin & Chino con el mapa básico de zhuyin \\
Wubi 86 & Chino mediante códigos de forma \\
US-International & Acentos latinos habituales \\
\end{choices}

Elija un resultado en la barra de candidatos. Seleccione
\key{Off — direct keyboard} en el mismo menú para volver al teclado estadounidense.
La selección afecta a todas las terminales y se guarda al salir de Inkline.

\topic{Tacto, lápiz y portapapeles}
\key{Dos dedos arriba/abajo:} recorren el historial. Arrastre hacia abajo para
ver salida antigua y hacia arriba para volver al indicador. Cada terminal
conserva 500 líneas físicas en RAM, además de la pantalla actual.

\key{Dos dedos izquierda/derecha:} cambian una posición por gesto.
Levante los dedos antes de volver a cambiar. Sepárelos o acérquelos para ajustar
el tamaño del texto.

\key{Un dedo:} arrastre para seleccionar. Alt/Opt derecha+C copia;
Alt/Opt derecha+V pega. Todas las terminales comparten un portapapeles en RAM
de hasta 128 KiB.

\key{Lápiz:} las aplicaciones que activan informes de ratón reciben pulsación,
movimiento y liberación como eventos del botón izquierdo, usando el protocolo
negociado, incluido SGR. En otros casos, arrastrar selecciona texto.
Mantenga \key{Shift} para forzar la selección local en aplicaciones con ratón.

\key{OSC 52} permite que programas locales o remotos por SSH lean o sustituyan
el portapapeles de Inkline. Es independiente del de los cuadernos.
Alt/Opt derecha+X copia e indica que la eliminación debe hacerse en el editor:
OSC 52 no puede borrar texto de un editor.

\key{Gráficos:} las transferencias kitty en línea y por memoria compartida
permanecen en RAM. Los modos de archivo y archivo temporal están desactivados.
Sixel funciona si se solicita explícitamente, pero no se anuncia automáticamente.
Los marcadores kitty, la reproducción de animaciones y algunos modos sixel
siguen incompletos.

\code{clear} y el borrado de pantalla completa eliminan texto y gráficos visibles.
Los gráficos fuera de pantalla se liberan cuando falta memoria. No hay caché
de imágenes en disco; las imágenes grandes siguen consumiendo la RAM limitada.

\topic{Teclado USB y máquina de escribir}
Abra \key{Settings → USB (página 3)}. \key{Network} restaura Ethernet USB;
\key{Send keyboard} convierte la tableta en teclado para otro ordenador;
\key{Receive keyboard} recibe un teclado por un adaptador OTG con alimentación.
Type Folio sigue disponible. Instale la utilidad Python de la web para enviar
pulsaciones; no hacen falta módulos Python adicionales.

Elija Send keyboard, conecte USB-C, seleccione el perfil y abra
\key{Open typewriter}. Enfoque un documento en el ordenador y toque \key{Start}.
Siempre se abre en pausa. \key{Escape} o \key{Stop} detiene el envío y libera las
teclas. Salir de esta pantalla, cambiar de terminal, perder el foco, abrir una
aplicación nativa o suspender también pausa. No hay reenvío automático al reconectar.

\begin{choices}
US / ASCII & Distribución US, Bloq Mayús desactivado. Rechaza texto no ASCII. \\
Mac & Fuente Unicode Hex Input, Bloq Mayús desactivado. Hasta U+FFFF;
rechaza caracteres de planos suplementarios. \\
Linux & Distribución US, Bloq Mayús desactivado. La aplicación debe aceptar
Ctrl+Shift+U, hexadecimal, Enter. No funciona en todas. \\
Windows & Instale WinCompose en el PC, Compose en Right Alt y entrada Unicode
activada. Distribución US, Bloq Mayús desactivado. Envía Compose, u, seis dígitos hexadecimales, Enter. \\
\end{choices}
WinCompose: \url{https://wincompose.info/}. Las reglas personalizadas no deben
sustituir esa secuencia. HID no tiene un canal universal de texto Unicode:
importan la configuración y las fuentes del receptor. Pruebe un texto corto.

\key{Input method} ofrece Off (US directo), Romaji, Pinyin, Zhuyin, Wubi 86
y US-International. Solo se envían caracteres confirmados. Cambiar de método pausa;
reanude con Start. \key{Off} devuelve la entrada US normal. \key{Unicode} o
Alt+Space abre el selector de caracteres para la salida USB.

Typewriter es un editor local cuyo documento permanece en RAM mientras Inkline
está abierto. La escritura normal modifica el documento y, con el modo en vivo,
envía cada carácter confirmado al ordenador. Copiar, cortar, pegar, seleccionar
y desplazar actúan en el editor local. \key{Files \& send} carga y guarda UTF-8
con cualquier nombre, permite 5--80 caracteres por segundo, envía todo el
documento y detiene el envío. Solo Save o la salida normal escriben el documento.
Un documento sin nombre usa \path{~/.inkline-typewriter-draft} al salir para su
recuperación. \key{Exit typewriter} vuelve a los ajustes USB.

\topic{Enviar archivos y restaurar la red}
Las órdenes \code{inkline-usb} y \code{keyboard-send} están en
\path{~/.local/bin}. \code{~/inkline usb}, \code{~/inkline type} e
\code{inkline-type} siguen como alias compatibles. Se conservan las herramientas Outpost.
\UsbExamples
Los archivos UTF-8 \key{no tienen un límite fijo de tamaño}. Se leen en bloques
pequeños, sin conservar el documento completo en RAM ni crear archivos temporales.
Los archivos que permiten rebobinar se validan enteros antes de enviar y se releen;
no los modifique durante el envío. Una tubería se valida al llegar: un error puede
ocurrir tras haber enviado parte del texto. Se quita el BOM UTF-8 inicial y
CRLF/CR pasan a saltos de línea. Enter y Tab actúan en la aplicación enfocada.

\code{--cps} (20 por defecto) o \code{--delay-ms} ajusta la velocidad; Unicode
requiere más informes. \code{--start-delay} deja tiempo para enfocar. Dry-run
muestra informes sin USB. Ctrl+C o \code{--stop} cancela y libera teclas.
Solo funciona un emisor a la vez. Si se llena la cola interactiva, se pausa.
Revise el documento antes de reintentar un envío parcial. El texto y el estado
temporal no se escriben en flash. El perfil se guarda con los demás ajustes al
salir normalmente.

\key{Network} o \code{inkline-usb network} restaura Ethernet. Salir de Inkline,
incluso tras un fallo o reinicio de emergencia, restaura la red. El modo USB no
se guarda como preferencia de arranque.
\UsbRecoveryExample
El temporizador es independiente del shell. Cambie de modo en la tableta o por
SSH Wi-Fi, nunca por SSH USB. Si falla, se intenta recuperar Ethernet. No se
quita el puerto a una sesión PPP satelital activa ni a un gadget desconocido;
también se aplica al salir. Opt+RightAlt+Backspace reinicia Inkline, pero cierra
todos sus terminales y pierde el trabajo sin guardar.

\topic{Abrir programas mediante atajos}
Registre atajos desde una shell de Inkline. Los argumentos son literales;
invoque una shell expresamente para usar su sintaxis. Estos programas leen
\code{~/.bashrc} mediante Bash.

Settings → Command shortcuts (página 2) permite asignar un comando a una letra. Elija Terminal o Native e-paper app y pulse Save. Remove pide confirmación; Reset draft descarta los cambios. T está bloqueada. Solo Save y la eliminación confirmada escriben en el almacenamiento. La herramienta de línea de comandos gestiona los mismos atajos.

Todos los atajos globales usan ahora Opt+RightAlt, incluidos T y la recuperación manteniendo Backspace. Use la tecla Opt independiente del Folio; en teclados USB, use Windows/Command (Super). Ctrl+Alt sigue disponible para Emacs. Cierre y vuelva a abrir Inkline tras actualizar; la instalación conserva los terminales abiertos.

\begin{shell}
~/inkline shortcut register e emacs
~/inkline shortcut register p python3
~/inkline shortcut list
~/inkline shortcut deregister p
\end{shell}

\key{Opt+RightAlt+E} abre Emacs en una posición libre.
\code{~/inkline run PROGRAM ARG...} inicia sin registrar un atajo.
Registrar de nuevo sustituye la asociación; eliminarla no cierra programas en
ejecución. Las teclas corresponden a posiciones físicas estadounidenses.
Entrecomille la puntuación y use rutas completas para programas propios,
que deben estar instalados.

Para una aplicación Qt/e-paper nativa que dibuja su propia pantalla:
\begin{shell}
~/inkline shortcut register a --epaper /home/root/my-app
\end{shell}
Solo se ejecuta una aplicación nativa a la vez, pausando Inkline y sus shells.
El lanzador configura Qt y cachés en RAM, descarta registros y reanuda Inkline
al terminar. Las aplicaciones deben usar entrada Qt sin capturar el teclado en exclusiva.

\begin{notice}
\key{Opt+RightAlt+T es permanente.} No se puede registrar ni eliminar.
Detiene la aplicación nativa y devuelve las terminales existentes.
No mantenga Ctrl para este atajo; la reasignación de Caps Lock solo actúa dentro de Inkline.
\end{notice}

\key{Emergencia:} mantenga Opt+RightAlt+Backspace dos segundos para detener la
aplicación y sus descendientes y reiniciar Inkline. Se cierran todas las
terminales y se pierde el trabajo sin guardar. Las aplicaciones root pueden
impedir la recuperación capturando entradas, modificando servicios o agotando
recursos. Reiniciar la tableta es el último recurso.

\topic{Instalar, actualizar y guardar la guía}
El paquete está destinado a \key{reMarkable 2 con firmware 3.27} y usa las
bibliotecas Qt 6.8 e-paper del sistema. Descargas y verificación de integridad:

\InstallLink

El instalador comprueba hardware, firmware, fuentes, arranque de Qt, directorios
temporales en RAM y swap desactivado. Las versiones residen en
\path{/home/root/.local/share/inkline}; el lanzador de teclado se inicia al arrancar.
La actualización conserva las terminales abiertas. \key{Cierre y vuelva a abrir
Inkline para usar la nueva versión.} Las versiones anteriores se guardan para
retroceder y ocupan espacio adicional.

Este único PDF contiene los nueve idiomas y viene en el paquete. Se intenta
importarlo a \key{My files} mediante la interfaz web USB cuando está activada
y los cuadernos están abiertos. Una importación correcta queda registrada
para no duplicar el mismo manual.

Si no funciona la importación automática, conecte USB, cierre Inkline y active
\key{USB web interface} en los ajustes de almacenamiento de los cuadernos.
Desde SSH, ejecute:
\begin{shell}
~/inkline manual
\end{shell}
También puede descargar el PDF e importarlo mediante el navegador USB o la
aplicación reMarkable de escritorio/móvil. El importador nunca modifica la
base de datos de cuadernos directamente ni interrumpe terminales. El manual
importado permanece en My files después de desinstalar Inkline.

Desde SSH, vuelva a los cuadernos o desinstale Inkline:
\begin{shell}
~/inkline stop
~/inkline uninstall
\end{shell}
La instalación de software opcional, incluido el perfil LaTeX recomendado,
se documenta en la web.

\topic{Funcionamiento, almacenamiento y RAM}
Qt entrega las teclas a un mapeador común para atajos y Caps Lock.
libghostty-vt codifica las teclas del terminal. Cada posición conecta su
programa mediante una pseudoterminal (PTY). Las shells reciben
\code{HOME=/home/root}; Bash interactivo lee \code{~/.bashrc}.

La salida pasa por el analizador del terminal y el decodificador sixel.
libghostty mantiene celdas, cursor, pantallas alternativas, historial e imágenes
kitty. El renderizador conserva una imagen gris y redibuja filas modificadas;
el backend e-paper de Qt actualiza el panel. Un demonio evdev independiente
gestiona atajos globales y el botón de encendido; systemd controla suspensión y reanudación.

Cada terminal limita las asignaciones del núcleo a 64 MiB, las imágenes kitty
a 16 MiB por pantalla y las imágenes sixel retenidas a 16 MiB. El renderizado
y las aplicaciones consumen memoria adicional. Las posiciones vacías no crean
terminales. Nueve aplicaciones grandes pueden superar la RAM disponible.

Gráficos temporales, cachés y sesiones usan RAM. El software instalado,
documentos, ajustes guardados al salir y archivos escritos por aplicaciones
usan almacenamiento persistente. Inkline no prohíbe que los programas escriban archivos.

El paquete Python opcional es Python estándar. Para reducir escrituras de
caché, añada a \code{~/.bashrc} en la tableta:
\begin{shell}
export PYTHONDONTWRITEBYTECODE=1
export PIP_NO_CACHE_DIR=1
\end{shell}
Ejecute \code{source ~/.bashrc} o abra otra shell. Use
\code{python3 -m pip install --no-compile PACKAGE} para evitar la compilación
de bytecode que pip realiza durante la instalación. El modo aislado de Python
ignora sus variables de entorno. La guía completa de Python está en la web.

Inkline usa GPL-3.0-or-later. Los paquetes incluyen los archivos fuente
correspondientes y avisos de dependencias. Los registros distinguen pruebas
automatizadas de comportamientos físicos pendientes de confirmación.
